\chapter{Temporal graphs}
\label{chap:model}

This chapter fixes the basic model: a temporal graph is a relation labeled by
discrete time, and we recall the handful of predicates on such a relation
(symmetry, looplessness, simplicity, being a clique) that the rest of the
library builds on. Formalized in \texttt{TemporalGraph/Basic.lean}.

\begin{definition}[Temporal graph]
  \label{def:temporal-graph}
  \lean{TemporalGraph}
  \leanok
  A \emph{temporal graph} on a vertex type $V$ is a relation
  $\mathrm{Adj} : V \to V \to \mathbb{N} \to \mathrm{Prop}$, where
  $\mathrm{Adj}(x,y,t)$ means there is an edge between $x$ and $y$ present at
  time $t$. No symmetry or irreflexivity is assumed a priori; see
  Definition~\ref{def:is-symm-loopless}.
\end{definition}

\begin{definition}[Symmetric, loopless]
  \label{def:is-symm-loopless}
  \uses{def:temporal-graph}
  \lean{TemporalGraph.IsSymm}
  \lean{TemporalGraph.IsLoopless}
  \leanok
  A temporal graph $G$ \emph{is symmetric} if $\mathrm{Adj}_G(x,y,t) \iff
  \mathrm{Adj}_G(y,x,t)$ for every $x,y,t$; it \emph{is loopless} if
  $\mathrm{Adj}_G(x,x,t)$ never holds.
\end{definition}

\begin{definition}[Simple, clique]
  \label{def:is-simple-clique}
  \uses{def:temporal-graph}
  \lean{TemporalGraph.IsSimple}
  \lean{TemporalGraph.IsClique}
  \leanok
  $G$ \emph{is simple} if between any two vertices there is at most one edge
  at a given time (i.e.\ $\mathrm{Adj}_G(x,y,n) \wedge \mathrm{Adj}_G(x,y,m)
  \Rightarrow n = m$). $G$ \emph{is a clique} if every pair of distinct
  vertices is joined by an edge at some time.
\end{definition}

\begin{definition}[Earliest / latest neighbor]
  \label{def:earliest-latest-neighbor}
  \uses{def:temporal-graph}
  \lean{TemporalGraph.EarliestEdge}
  \lean{TemporalGraph.LatestEdge}
  \lean{TemporalGraph.EarliestNeighbor}
  \lean{TemporalGraph.LatestNeighbor}
  \leanok
  For $x \in V$, an edge $(x,y,n)$ is the \emph{earliest edge} of $x$ towards
  $y$ if $n$ is minimal among times at which $x$ has an edge; it is the
  \emph{latest edge} if $n$ is maximal. $y$ is an \emph{earliest} (resp.\
  \emph{latest}) \emph{neighbor} of $x$ if such an edge exists. These notions
  drive the dismountability argument of Chapter~\ref{chap:dismountability}.
\end{definition}

\begin{definition}[Spanning sub-temporal-graph order]
  \label{def:spanning-order}
  \uses{def:temporal-graph}
  \lean{TemporalGraph.IsSpanningSubgraph}
  \leanok
  $G$ is a \emph{spanning sub-temporal-graph} of $H$ (written $G \le H$) if
  every edge of $G$ is an edge of $H$: $\mathrm{Adj}_G(x,y,n) \Rightarrow
  \mathrm{Adj}_H(x,y,n)$ for all $x,y,n$.
\end{definition}

\begin{definition}[General sub-temporal-graph]
  \label{def:is-subgraph-of}
  \uses{def:temporal-graph}
  \lean{TemporalGraph.IsSubgraphOf}
  \leanok
  Let $P, P'$ be predicates on $V$. A temporal graph $G$ carried by the
  vertices satisfying $P$ is a \emph{sub-temporal-graph} of $H$, carried by
  those satisfying $P'$, when $P$ implies $P'$ and every edge of $G$ is an
  edge of $H$ between the corresponding vertices. This is the general
  inclusion: it compares graphs living on two different vertex subsets, and
  Definition~\ref{def:spanning-order} is its special case $P = P'$.
\end{definition}

\begin{definition}[Arcs and edges]
  \label{def:arcset-edgeset}
  \uses{def:temporal-graph,def:is-symm-loopless}
  \lean{TemporalGraph.arcSet}
  \lean{TemporalGraph.edgeSet}
  \lean{TemporalGraph.mem_edgeSet_iff_or}
  \leanok
  The \emph{arcs} of $G$ are the triples $(x,y,n)$ with $\mathrm{Adj}_G(x,y,n)$:
  an oriented notion, counted by $\mathrm{arcSize}$. The \emph{edges} of $G$
  are unordered: pairs $(\{x,y\}, n)$ such that $\mathrm{Adj}_G(x,y,n)$ or
  $\mathrm{Adj}_G(y,x,n)$, counted by $\mathrm{size}$. An unordered pair is
  therefore an edge exactly when one of its orientations is an arc, with no
  hypothesis; when $G$ is symmetric the two orientations coincide and
  membership reduces to $\mathrm{Adj}_G(x,y,n)$.
\end{definition}

\begin{definition}[Symmetrization]
  \label{def:symmetrize}
  \uses{def:temporal-graph,def:is-symm-loopless,def:spanning-order}
  \lean{TemporalGraph.symmetrize}
  \lean{TemporalGraph.le_symmetrize}
  \lean{TemporalGraph.symmetrize_le}
  \leanok
  The \emph{symmetrization} of $G$ has $\mathrm{Adj}(x,y,n)$ whenever
  $\mathrm{Adj}_G(x,y,n)$ or $\mathrm{Adj}_G(y,x,n)$. It is symmetric, it lies
  above $G$, and it lies below every symmetric graph above $G$ — that is, it is
  the least symmetric temporal graph above $G$. It adds no loop, so it
  preserves looplessness.
\end{definition}

\begin{definition}[Building a temporal graph]
  \label{def:from-arc-edge-set}
  \uses{def:arcset-edgeset}
  \lean{TemporalGraph.fromArcSet}
  \lean{TemporalGraph.fromEdgeSet}
  \leanok
  A set of arcs yields a temporal graph directly, with no symmetry assumed. A
  set of unordered edges yields one that is symmetric by construction, each
  pair being read in both directions; its edge set is the set one started
  from.
\end{definition}
